\section*{الفصل \ref{ch:g5:series-parallel-circuits} --- الدارات على التسلسل وعلى التفرع}

\begin{solution}{exo:g5:series-parallel-circuits:1}
\omterm{def:g5:series-parallel-circuits:parallel}{الفرع} أحد الطرق المنفصلة التي تتشعّب إليها الدارة؛ \omterm{def:g5:series-parallel-circuits:parallel}{والعقدة} نقطة
تتشعّب عندها الطرق أو تلتقي. ومصباحان \omterm{def:g5:series-parallel-circuits:parallel}{على التفرع} يركبان فرعين
حاملين للمصابيح.
\end{solution}

\begin{solution}{exo:g5:series-parallel-circuits:2}
الاستقلال (فالفراغ لا يوقف إلا فرعه)؛ والسطوع الكامل (فيسطع كل
\omterm{def:g3:first-electric-circuit:bulb}{مصباح} كأنه وحده، \omterm{def:g3:first-electric-circuit:battery}{والبطارية} تنفد أسرع في المقابل)؛ ومواقع الأمر
(فقاطعة \omterm{def:g5:series-parallel-circuits:parallel}{الفرع} تحكم فرعها، \omterm{ex:g3:first-electric-circuit:switch}{والقاطعة} قبل التشعّب تحكم الجميع).
\end{solution}

\begin{solution}{exo:g5:series-parallel-circuits:3}
يظلّ \omterm{def:g3:first-electric-circuit:bulb}{المصباح} 2 مشتعلًا غير مضطرب --- فطريقه الخاص كامل. أما \omterm{def:g4:series-circuits:series}{على
التسلسل} فقد أظلم، لأن المصباحين هناك كانا يتقاسمان طريقًا واحدًا،
وأيّ فراغ يقطعه على كليهما.
\end{solution}

\begin{solution}{exo:g5:series-parallel-circuits:4}
الزوج B (\omterm{def:g5:series-parallel-circuits:parallel}{على التفرع}) أسطع في كل \omterm{def:g3:first-electric-circuit:bulb}{مصباح} --- فكل واحد بسطوع كامل.
والزوج B أيضًا ينفد بطاريته أسرع: فهو يغذّي فرعين كاملين دفعة
واحدة.
\end{solution}

\begin{solution}{exo:g5:series-parallel-circuits:5}
القاطعتان مغلقتان: يسطع المصباحان. والرئيسية مفتوحة: يعتم
الاثنان --- فهي تحكم الطريق المشترك. والرئيسية مغلقة وقاطعة
\omterm{def:g5:series-parallel-circuits:parallel}{الفرع} مفتوحة: يسطع \omterm{def:g3:first-electric-circuit:bulb}{المصباح} 1 ويعتم \omterm{def:g3:first-electric-circuit:bulb}{المصباح} 2.
\end{solution}

\begin{solution}{exo:g5:series-parallel-circuits:6}
\omterm{def:g4:series-circuits:series}{على التسلسل} يُظلم \omterm{def:g3:first-electric-circuit:bulb}{مصباح} ميت واحد (أو \omterm{def:g3:first-electric-circuit:bulb}{مصباح} مطفأ!) البيتَ كله،
ويخفت كل شيء كلما انضمّ \omterm{def:g1:light-and-shadows:source}{ضوء} جديد. والراحات التي يمنحها التفرع:
أن تحكم قاطعة كل غرفة غرفتَها وحدها؛ وأن يبقى سائر البيت مضيئًا
مع احتراق \omterm{def:g3:first-electric-circuit:bulb}{مصباح} المطبخ؛ وأن يشتعل كل \omterm{def:g3:first-electric-circuit:bulb}{مصباح} بسطوع كامل.
\end{solution}

\begin{solution}{exo:g5:series-parallel-circuits:7}
قبل التشعّب --- على الطريق المشترك الذي يدخل منه التيار إلى
البيت: أي \omterm{ex:g3:first-electric-circuit:switch}{القاطعة} الرئيسية في علبة المصاهر. وهي تقطع البيت كله
دفعة واحدة، وتُستعمل للإصلاح والطوارئ.
\end{solution}

\begin{solution}{exo:g5:series-parallel-circuits:8}
\omterm{def:g5:series-parallel-circuits:parallel}{على التفرع}: فإطفاء \omterm{def:g1:light-and-shadows:source}{ضوء} السقف لا يقتل \omterm{def:g3:first-electric-circuit:bulb}{مصباح} السرير، ونزع قابس
\omterm{def:g3:first-electric-circuit:bulb}{المصباح} لا يُظلم السقف --- فكل واحد يركب فرعه الخاص.
\end{solution}

\begin{solution}{exo:g5:series-parallel-circuits:9}
مع قواطعنا البسيطة تظلّ التصاميم تفشل: \omterm{ex:g3:first-electric-circuit:switch}{فالقاطعة} \omterm{def:g4:series-circuits:series}{على التسلسل} مع
\omterm{def:g3:first-electric-circuit:bulb}{المصباح} لا تُفتح إلا من طرف واحد؛ وقاطعتان بسيطتان \omterm{def:g4:series-circuits:series}{على التسلسل}
تقتضيان أن تكونا مغلقتين \emph{معًا} (فيستطيع كل طرف أن يُطفئ
\omterm{def:g1:light-and-shadows:source}{الضوء} لا أن يُشعله دائمًا)؛ وقاطعتان \omterm{def:g5:series-parallel-circuits:parallel}{على التفرع} تقتضيان أن
تكونا مفتوحتين معًا للعتمة. فراحة الممرّ تحتاج إلى قاطعة
\emph{تختار بين طريقين} بدل أن تقطع طريقًا واحدًا --- وهي قاطعة
التبديل، وتستعملها الممرّات الحقيقية أزواجًا متوافقة.
\end{solution}

\begin{solution}{exo:g5:series-parallel-circuits:10}
أما الراحة: فالمصابيح تركب \omterm{def:g5:series-parallel-circuits:parallel}{فروعًا} متفرّعة، فلا يُظلم \omterm{def:g3:first-electric-circuit:bulb}{المصباح}
الميت إلا نفسه. وأما التحذير: فكل \omterm{def:g3:first-electric-circuit:bulb}{مصباح} ناقص ينزع فرعًا، فتدفع
\omterm{def:g3:first-electric-circuit:battery}{البطارية} (أو مغذّي السلسلة الصغير) كامل جهدها في الفروع الباقية
--- وإذا كثر الناقص نال الباقون أكثر ممّا بُنوا له، فسخنوا أكثر
وماتوا صغارًا.
\end{solution}

\begin{solution}{exo:g5:series-parallel-circuits:11}
على امتداد الطريق الخاصّ بقاطعة \omterm{def:g5:series-parallel-circuits:parallel}{الفرع}: بين قاطعة \omterm{def:g5:series-parallel-circuits:parallel}{الفرع} \omterm{def:g3:first-electric-circuit:bulb}{والمصباح}
2 (قبل أن يعود فرعهما إلى سائر الدارة). فهناك يضيء بالضبط حين
يعمل ذلك \omterm{def:g5:series-parallel-circuits:parallel}{الفرع} --- ويكون \omterm{def:g4:series-circuits:series}{على التسلسل} مع \omterm{def:g3:first-electric-circuit:bulb}{المصباح} 2، إذ يركبان
\omterm{def:g5:series-parallel-circuits:parallel}{الفرع} نفسه.
\end{solution}
